%%
%% Neighbourhood Structures and Ranking Operators
%% for the Interval Job Shop Scheduling Problem
%%
%% Built on Elsevier's elsarticle-template-num.tex
%% (numbered references; preprint single-column layout).
%%
%% Compile with:  latexmk -pdf main
%%

%% Submission style (what Elsevier asks for at submission time):
\documentclass[preprint,12pt]{elsarticle}

%% Journal-style previews (only for visualisation; do NOT submit in these):
%% \documentclass[final,1p,times]{elsarticle}                   % 1 col journal
%% \documentclass[final,3p,times,twocolumn]{elsarticle}         % 2 col, narrow
%% \documentclass[final,5p,times,twocolumn]{elsarticle}         % 2 col, wide

%% The amssymb package provides various useful mathematical symbols
\usepackage{amssymb}
%% The amsmath package provides various useful equation environments.
\usepackage{amsmath}

%% Theorem-like environments via base LaTeX (no amsthm needed)
\newtheorem{theorem}{Theorem}
\newtheorem{proposition}{Proposition}
\newtheorem{corollary}{Corollary}
\newtheorem{definition}{Definition}

%% Proof environment
\newenvironment{proof}%
  {\noindent\textit{Proof.}\enspace}%
  {\hfill$\square$\par\medskip}

%% Show a placeholder box if a figure file is not yet on disk.
\newcommand{\figorplaceholder}[3]{%
  \IfFileExists{#1.pdf}{%
    \includegraphics[width=#2]{#1}%
  }{%
    \fbox{\parbox[c][#3][c]{#2}{\centering\itshape%
      Placeholder for figure \texttt{\detokenize{#1}.pdf}\\(file not yet on disk)}}%
  }%
}

\graphicspath{{../../figures/}{./figures/}}

\journal{Computers \& Operations Research}

\begin{document}

\begin{frontmatter}

%% Title, authors and addresses

\title{Neighbourhood Structures and Ranking Operators for the Interval Job Shop Scheduling Problem}

\author[uo]{Hern\'an D\'{\i}az Rodr\'{\i}guez\corref{cor1}}
\ead{diazhernan@uniovi.es}
\cortext[cor1]{Corresponding author}

\affiliation[uo]{organization={Department of Computing, University of Oviedo},
            addressline={Campus de Gij\'on},
            city={Gij\'on},
            postcode={33204},
            country={Spain}}

%% Abstract
\begin{abstract}
The Job Shop Scheduling Problem with interval-valued processing times (IJSP) models real-world uncertainty where durations are unknown but bounded. Comparing schedules then requires an interval ranking operator, and the classical notion of critical path must be extended to both extreme scenarios. I formalise \emph{extreme critical paths} --- paths critical in the best-case or worst-case scenario graph --- and use them to define a neighbourhood $H(\sigma)$ for local search. I prove that $H(\sigma)$ is feasible, connected and loses no improving neighbours under any admissible interval ranking. Five neighbourhood variants ($N_1$, $N_2$, $N_3$, $N_{\rm ext}$, $N_8$) are formulated within this framework and combined with four ranking operators (EV, LEX1, LEX2, YX) on 82 benchmark instances of up to 1000 operations. Two findings are noteworthy: \emph{(i)} after irace-based tuning, $N_{\rm ext}$ is the strongest neighbourhood, followed by $N_1$ and $N_2$ (negligible to small effects), with $N_8$ and $N_3$ trailing by a medium effect; and \emph{(ii)} LEX2 systematically yields makespan intervals 5--11\% narrower than the alternatives. A float analysis of both extreme graphs explains this effect mechanistically: LEX2 directly tightens the worst-case critical path rather than leaving structural slack. A validation on the 82 deterministic counterparts (CRISP) confirms that the relative ranking of neighbourhoods carries over to the classical setting, and that $N_8$ becomes competitive when uncertainty disappears.
\end{abstract}

%% Graphical abstract
\begin{graphicalabstract}
%\includegraphics[width=\linewidth]{graphical_abstract}
\end{graphicalabstract}

%% Research highlights
\begin{highlights}
\item Unified extreme-critical-path framework for neighbourhoods in the interval JSP.
\item Feasibility, connectivity and no-loss proofs for any admissible interval ranking.
\item LEX2 yields 5--11\% narrower makespan intervals than EV, YX and LEX1.
\item Float analysis: LEX2 directly tightens the worst-case critical path.
\item Slack-filtered $N_8$ is CRISP-good but IJSP-bad: rejects disagreement moves.
\end{highlights}

%% Keywords
\begin{keyword}
Job shop scheduling \sep interval uncertainty \sep critical paths \sep neighbourhood structures \sep interval ranking \sep local search
\end{keyword}

\end{frontmatter}

%% main text

\section{Introduction}
\label{sec1}

The Job Shop Scheduling Problem (JSP) is one of the most thoroughly studied combinatorial optimisation problems in operations research. In many real-world applications, however, processing times are not known precisely at the time of planning: they depend on operator experience, machine wear, material quality, and other variable factors. A natural and tractable model for such uncertainty is to represent each processing time as a \emph{closed interval} $[p^-, p^+]$ where $p^-$ and $p^+$ are the minimum and maximum possible durations \cite{Moore2009,GonzalezRodriguez2008,Lei2012c}.

The resulting \emph{Interval Job Shop Scheduling Problem (IJSP)} requires comparing schedules whose makespan is itself an interval, making the notion of optimality dependent on the choice of \emph{interval ranking operator} \cite{Bustince2013,Karmakar2012}. Furthermore, the classical concept of a critical path --- fundamental to neighbourhood design in the deterministic JSP \cite{vanLaarhoven1992,Nowicki1996,Nowicki2005} --- must be reconsidered: which path is critical when processing times are uncertain?

\subsection{Contributions}
\label{sec1.1}

This paper makes the following contributions.
\begin{enumerate}
\item \textbf{Formal definition of extreme criticality for IJSP} (Section~\ref{sec3}). I define a task or arc as \emph{extreme critical} if it lies on a critical path of the best-case graph $G^-(\sigma)$ or the worst-case graph $G^+(\sigma)$, and prove that this definition encompasses all necessarily critical tasks and is contained within the set of possibly critical tasks, making it a computationally tractable and theoretically sound approximation of the general criticality concept of \cite{Artigues2015}.

\item \textbf{Neighbourhood $H(\sigma)$ based on extreme critical arcs} (Section~\ref{sec4}). I define $H(\sigma)$ as the set of all feasible processing orders reachable by reversing a single extreme critical disjunctive arc, and prove three structural properties: feasibility (every move yields a feasible order), no loss of improving neighbours (moves outside $H(\sigma)$ cannot improve the makespan under any admissible ranking), and connectivity (every non-optimal order can reach a global optimum via a finite sequence of moves in $H(\sigma)$).

\item \textbf{Five neighbourhood variants within the $H(\sigma)$ framework} (Section~\ref{sec5}). I formulate $N_1$, $N_2$, $N_3$, $N_{\rm ext}$ and a new variant $N_8$ (boundary-only with extra slack-based filtering) as instances of $H(\sigma)$, characterising their theoretical size and the cost of generating them.

\item \textbf{Empirical study on 82 benchmark instances} (Section~\ref{sec6}). Three experimental phases are reported: \emph{(a)} a $4 \times 4$ comparison of operators and neighbourhoods at a common parameter setting (Section~\ref{sec6.2}); \emph{(b)} a head-to-head comparison of the five neighbourhoods after irace-based hyperparameter tuning, isolating the structural effect of the neighbourhood (Section~\ref{sec6.5}); and \emph{(c)} a validation phase on the 82 deterministic counterparts of the same instances (CRISP, Section~\ref{sec6.6}).

\item \textbf{Discovery and structural explanation of the LEX2--width effect} (Sections~\ref{sec6.3} and~\ref{sec6.4}). LEX2 systematically produces solutions with $\sim$5--11\% narrower makespan intervals than the other operators. A float analysis of both extreme graphs falsifies the natural ``more slack'' hypothesis and shows that the effect arises mechanistically: LEX2 tightens the $G^+(\sigma)$ critical path directly.
\end{enumerate}

\subsection{Paper Organisation}
\label{sec1.2}
Section~\ref{sec2} introduces the IJSP formally and reviews related work on critical paths and rankings. Section~\ref{sec3} defines extreme criticality. Section~\ref{sec4} defines the neighbourhood $H(\sigma)$ and proves its three structural properties. Section~\ref{sec5} describes the five neighbourhood variants. Section~\ref{sec6} reports the experimental study. Section~\ref{sec7} discusses the findings, and Section~\ref{sec8} concludes.


\section{Background and Related Work}
\label{sec2}

\subsection{The Deterministic Job Shop Problem}
\label{sec2.1}

The JSP consists of $n$ jobs $\{J_1,\ldots,J_n\}$ and $m$ machines $\{M_1,\ldots,M_m\}$. Each job $J_i$ is a sequence of operations $O_{i,1},\ldots,O_{i,m}$, where operation $O_{i,j}$ must be processed on machine $M_j$ for exactly $p_{i,j}$ time units. Operations within a job must respect their precedence order, and each machine processes at most one operation at a time. The goal is to find a \emph{feasible processing order} $\sigma$ minimising the makespan $C_{\max}$.

The \emph{solution graph} $G(\sigma) = (V, A \cup D(\sigma))$ represents a feasible order: $A$ contains precedence arcs within jobs, and $D(\sigma)$ contains disjunctive arcs imposed by the order $\sigma$ on each machine. The makespan equals the length of the longest path in $G(\sigma)$. The classical insight of \cite{vanLaarhoven1992} is that the neighbourhood for local search should restrict to reversals of arcs on critical paths, ensuring feasibility, while \cite{Nowicki1996} shows that reversals of \emph{interior} arcs of a critical block cannot improve the makespan, allowing further pruning.

\subsection{Interval Arithmetic and Ranking Operators}
\label{sec2.2}

When processing times are uncertain, we represent each as $\mathbf{p}=[p^-, p^+] \in \mathbb{IR}$, where $\mathbb{IR}$ denotes the set of closed intervals. The IJSP requires three arithmetic operations on $\mathbb{IR}$ \cite{Moore2009}:
\begin{align}
\mathbf{a} + \mathbf{b}        &= [a^- + b^-,\, a^+ + b^+], \\
\max(\mathbf{a},\mathbf{b})    &= [\max(a^-,b^-),\, \max(a^+,b^+)].
\end{align}

Since there is no natural total order on $\mathbb{IR}$, comparing schedules requires choosing a \emph{ranking operator}. We consider four rankings widely used in the literature \cite{Bustince2013,Xu2006,Lei2012c}:
\begin{align}
\mathbf{a} \leq_{\rm LEX1} \mathbf{b} &\iff a^- < b^- \,\vee\, (a^- = b^- \wedge a^+ \leq b^+) \\
\mathbf{a} \leq_{\rm LEX2} \mathbf{b} &\iff a^+ < b^+ \,\vee\, (a^+ = b^+ \wedge a^- \leq b^-) \\
\mathbf{a} \leq_{\rm YX}   \mathbf{b} &\iff (a^- + a^+) < (b^- + b^+) \,\vee\, ((a^- + a^+) = (b^- + b^+) \wedge (a^+ - a^-) \leq (b^+ - b^-)) \\
\mathbf{a} \leq_{\rm EV}   \mathbf{b} &\iff \tfrac{a^- + a^+}{2} \leq \tfrac{b^- + b^+}{2}
\end{align}
LEX1, LEX2 and YX are \emph{admissible} total orders on $\mathbb{IR}$ \cite{Bustince2013}: they refine the partial product order $\leq_{2}$ defined by $\mathbf{a} \leq_{2} \mathbf{b} \iff a^- \leq b^- \wedge a^+ \leq b^+$. EV (midpoint) is not strictly a total order, but also refines $\leq_{2}$. This refinement property is essential for the theoretical guarantees of Section~\ref{sec4}: any move that simultaneously decreases both $C^-_{\max}$ and $C^+_{\max}$ is an improvement under all four rankings considered here.

\subsection{Related Work on IJSP Neighbourhoods}
\label{sec2.3}

The first neighbourhood for the fuzzy/interval JSP was proposed in \cite{GonzalezRodriguez2008}, extending the van Laarhoven neighbourhood \cite{vanLaarhoven1992} to use possibly critical arcs. This corresponds to what we call $N_1$ in our framework. Subsequent work introduced $N_2$ as a filtered version of $N_1$ that retains only block-boundary arcs \cite{N2fjsp}, reducing neighbourhood size by $\sim$64\% with statistically identical solution quality. $N_3$ extends the neighbourhood to consider subsequence reversals \cite{N3fjsp}.

The Nowicki--Smutnicki extended neighbourhood \cite{Nowicki2005}, which we adapt here under the name $N_{\rm ext}$, was originally proposed for the crisp JSP and provides strong empirical performance. Its adaptation to the IJSP, presented in this work, takes advantage of the fact that the boundary-arc result of \cite{Nowicki1996} can be lifted simultaneously to both extreme graphs $G^-(\sigma)$ and $G^+(\sigma)$. We also propose a fifth variant, $N_8$, which uses a stricter slack-based filter on top of $N_2$.


\section{Extreme Criticality for IJSP}
\label{sec3}

\subsection{Configurations and Extreme Scenario Graphs}
\label{sec3.1}

A \emph{configuration} $\omega = (p_1,\ldots,p_o)$ assigns a deterministic duration $p_i \in \mathbf{p}_i$ to each operation~$i$; let $\Omega$ denote the set of all configurations. Given a feasible processing order $\sigma$, two \emph{extreme scenario graphs} are of particular interest:
\begin{itemize}
\item $G^-(\sigma)$: identical to $G(\sigma)$ but with each arc $(x,y)$ weighted by $p_{x}^{-}$ (best case).
\item $G^+(\sigma)$: identical to $G(\sigma)$ but with each arc $(x,y)$ weighted by $p_{x}^{+}$ (worst case).
\end{itemize}
These correspond to the extreme configurations $\omega^{+}(\emptyset)$ and $\omega^{+}(V)$ of \cite{Artigues2015}. The makespan interval satisfies
\begin{equation}
\mathbf{C}_{\max}(\sigma) = [C^-_{\max}(\sigma),\, C^+_{\max}(\sigma)],
\end{equation}
where $C^-_{\max}(\sigma)$ is the length of the longest path in $G^-(\sigma)$ and $C^+_{\max}(\sigma)$ is the length of the longest path in $G^+(\sigma)$. An illustrative example is shown in Figure~\ref{fig1}: two jobs ($J_1$: $O_{11}(M_1,[2,6])\!\to\!O_{12}(M_2,[2,3])$; $J_2$: $O_{21}(M_2,[4,4])\!\to\!O_{22}(M_1,[1,4])$), with schedule $\sigma$ fixing $M_1$ order $(O_{11},O_{22})$ and $M_2$ order $(O_{21},O_{12})$. In $G^-(\sigma)$ the critical path runs through the $M_2$ disjunctive arc ($C_{\max}^-=6$); in $G^+(\sigma)$ it runs through the $M_1$ disjunctive arc ($C_{\max}^+=10$). The critical arcs in the two extreme graphs differ, motivating the need to consider both.

\begin{figure}[t]
\centering
\figorplaceholder{fig1_example_graphs}{0.9\linewidth}{6cm}
\caption{Illustrative example. (a) Best-case graph $G^-(\sigma)$ with its critical path highlighted in red. (b) Worst-case graph $G^+(\sigma)$ with its (different) critical path highlighted in blue. Precedence arcs are drawn solid, disjunctive arcs dashed.}
\label{fig1}
\end{figure}

\subsection{Heads, Tails and Floats}
\label{sec3.2}

For each extreme graph $G^{*}(\sigma)$ ($* \in \{-,+\}$), the \emph{head} $r_{x}^{*}$ of task $x$ is the length of the longest path from the dummy start node to $x$, and the \emph{tail} $q_{x}^{*}$ is the length of the longest path from the completion of $x$ to the end node. The \emph{float} of $x$ in $G^{*}(\sigma)$ is
\begin{equation}
F_{x}^{*}(\sigma) = C_{\max}^{*}(\sigma) - r_{x}^{*} - p_{x}^{*} - q_{x}^{*}.
\end{equation}
A task is critical in $G^{*}(\sigma)$ iff $F_{x}^{*}(\sigma) = 0$.

\subsection{Definition and Connection to Possibly/Necessarily Critical Tasks}
\label{sec3.3}

\begin{definition}[Extreme critical path]
\label{def1}
A path $P$ in $G(\sigma)$ is \emph{extreme critical} if it is critical in $G^-(\sigma)$ or in $G^+(\sigma)$. A node, arc or block is extreme critical if it lies on an extreme critical path.
\end{definition}

This definition connects to the general criticality concepts of \cite{Artigues2015} as follows.

\begin{proposition}
\label{prop1}
Every extreme critical task is possibly critical. Every necessarily critical task is extreme critical.
\end{proposition}

\begin{proof}
\textit{(i) Extreme critical $\Rightarrow$ possibly critical.}
Let task $x$ be extreme critical; then $F_x^-(\sigma)=0$ or $F_x^+(\sigma)=0$.
The extreme configurations $\omega^- = (p_i^-)_{i}$ and $\omega^+ = (p_i^+)_{i}$ both belong to $\Omega$
(since $p_i^{\pm}\in[p_i^-,p_i^+]$ by definition).
In case $F_x^-(\sigma)=0$: the float of $x$ under configuration $\omega^-$ equals
$F_x(\sigma,\omega^-)=C_{\max}^-(\sigma)-r_x^- - p_x^- - q_x^- = F_x^-(\sigma)=0$,
so $x$ lies on a critical path for $\omega^-\in\Omega$, making it possibly critical \cite{Artigues2015}.
The case $F_x^+(\sigma)=0$ is analogous with $\omega^+$.

\textit{(ii) Necessarily critical $\Rightarrow$ extreme critical.}
If $x$ is necessarily critical, then $F_x(\sigma,\omega)=0$ for every $\omega\in\Omega$.
Taking $\omega=\omega^-\in\Omega$ gives $F_x^-(\sigma)=0$, so $x$ is critical in $G^-(\sigma)$, hence extreme critical.
\end{proof}

\medskip

The practical consequence is that extreme criticality offers a computationally tractable approximation to the exact criticality predicates: computing critical paths in $G^-(\sigma)$ and $G^+(\sigma)$ each costs $O(|E|)$ via standard longest-path algorithms, while computing the full set of possibly critical tasks is NP-hard in general \cite{Artigues2015}.


\section{The Neighbourhood $H(\sigma)$ and Its Properties}
\label{sec4}

\subsection{Definition}
\label{sec4.1}

\begin{definition}[Neighbourhood $H(\sigma)$]
\label{def2}
Let $\sigma_{(v)}$ denote the processing order obtained from $\sigma$ by reversing a single disjunctive arc $v \in D(\sigma)$. The neighbourhood is
\begin{equation}
H(\sigma) = \{ \sigma_{(v)} \,:\, v \in D(\sigma) \text{ is extreme critical}\}.
\end{equation}
For comparison, the corresponding neighbourhood based on \emph{possibly} critical arcs is
\begin{equation}
H'(\sigma) = \{ \sigma_{(v)} \,:\, v \in D(\sigma) \text{ is possibly critical}\}.
\end{equation}
By Proposition~\ref{prop1}, $H(\sigma) \subseteq H'(\sigma)$, with strict inclusion in general.
\end{definition}

\subsection{No Loss of Improving Neighbours}
\label{sec4.2}

\begin{proposition}
\label{prop2}
Let $\sigma \in \Sigma$ be a feasible order and let $v \in D(\sigma)$ be a disjunctive arc that is \emph{not} extreme critical in $G(\sigma)$. Then
\begin{equation}
C^-_{\max}(\sigma) \leq C^-_{\max}(\sigma_{(v)}) \quad\text{and}\quad C^+_{\max}(\sigma) \leq C^+_{\max}(\sigma_{(v)}),
\end{equation}
and consequently $\mathbf{C}_{\max}(\sigma) \leq_{R} \mathbf{C}_{\max}(\sigma_{(v)})$ for every ranking $R \in \{{\rm EV}, {\rm LEX1}, {\rm LEX2}, {\rm YX}\}$.
\end{proposition}

\begin{proof}
We prove $C_{\max}^*(\sigma)\le C_{\max}^*(\sigma_{(v)})$ for $*\in\{-,+\}$; both cases are identical.

Since $(x,y)$ is not extreme critical, it lies on no critical path of $G^*(\sigma)$.
Let $P^*$ be any critical path of $G^*(\sigma)$, of length $C_{\max}^*(\sigma)$.
$P^*$ does not contain arc $(x,y)$.
Since $G^*(\sigma_{(v)})$ differs from $G^*(\sigma)$ only by replacing $(x,y)$ with $(y,x)$,
every arc of $P^*$ is still present in $G^*(\sigma_{(v)})$.
Hence $P^*$ is a valid directed path in $G^*(\sigma_{(v)})$ of weight $C_{\max}^*(\sigma)$,
so $C_{\max}^*(\sigma_{(v)})\ge C_{\max}^*(\sigma)$.

Applied to both extreme graphs: $C_{\max}^-(\sigma_{(v)})\ge C_{\max}^-(\sigma)$ and
$C_{\max}^+(\sigma_{(v)})\ge C_{\max}^+(\sigma)$,
i.e., $\mathbf{C}_{\max}(\sigma)\le_2\mathbf{C}_{\max}(\sigma_{(v)})$.
Since every ranking $R\in\{\mathrm{EV},\mathrm{LEX1},\mathrm{LEX2},\mathrm{YX}\}$ refines $\le_2$,
we obtain $\mathbf{C}_{\max}(\sigma)\le_R\mathbf{C}_{\max}(\sigma_{(v)})$.
\end{proof}

\medskip

\begin{corollary}
\label{cor1}
Neighbours in $H'(\sigma) \setminus H(\sigma)$ can never improve the makespan under any of the four rankings considered.
\end{corollary}

Corollary~\ref{cor1} justifies using $H(\sigma)$ instead of the larger $H'(\sigma)$: the reduction is gained at zero cost in solution quality.

\subsection{Feasibility}
\label{sec4.3}

\begin{theorem}[Feasibility]
\label{thm1}
For every feasible $\sigma \in \Sigma$, the reversal of any extreme critical arc $v = (x,y) \in D(\sigma)$ produces a feasible processing order $\sigma_{(v)} \in \Sigma$. Hence $H(\sigma) \subseteq \Sigma$.
\end{theorem}

\begin{proof}
Suppose for contradiction that $G(\sigma_{(v)})$ contains a directed cycle.
Since $G(\sigma)$ is acyclic and $G(\sigma_{(v)})=(G(\sigma)\setminus\{(x,y)\})\cup\{(y,x)\}$,
any cycle in $G(\sigma_{(v)})$ must traverse the newly added arc $(y,x)$.
Therefore there exists a directed path $\pi$ from $x$ to $y$ in $G(\sigma_{(v)})$ that avoids $(y,x)$.
All arcs of $\pi$ belong to $G(\sigma)\setminus\{(x,y)\}$ (the arc $(x,y)$ was removed, and $\pi$ uses none of the new arcs); so $\pi$ is a valid path from $x$ to $y$ in $G(\sigma)$ not using $(x,y)$.

In the standard JSP, $x$ and $y$ belong to the same machine but to different jobs; hence no job-precedence arc of $A$ runs directly from $x$ to $y$.
Therefore $\pi$ has at least one intermediate operation $u$ distinct from $x$ and $y$.
The length of $\pi$ in $G^*(\sigma)$ is
\[
  \ell(\pi) \;=\; p_x^* + p_u^* + \cdots \;\ge\; p_x^* + p_u^* \;>\; p_x^*,
\]
since all processing times satisfy $p_z^*\ge p_z^->0$.

On the other hand, $(x,y)\in CP^*(\sigma)$ implies that the head of $y$ is tight through the arc $(x,y)$:
\[
  r_y^* \;=\; r_x^* + p_x^*.
\]
But the path $\pi$ from $x$ to $y$ in $G(\sigma)$ yields
\[
  r_y^* \;\ge\; r_x^* + \ell(\pi) \;>\; r_x^*+p_x^* \;=\; r_y^*,
\]
a contradiction.
Hence $G(\sigma_{(v)})$ is acyclic, and $\sigma_{(v)}\in\Sigma$.
\end{proof}

\subsection{Connectivity}
\label{sec4.4}

\begin{theorem}[Connectivity]
\label{thm2}
For every non-optimal $\sigma \in \Sigma$, there exists a finite sequence of moves in $H(\sigma)$ leading to a globally optimal order $\sigma^{\star}$.
\end{theorem}

\begin{proof}
Fix any globally optimal order $\sigma^*\in\Sigma^*$ that minimises the \emph{inversion count}
\[
  M(\sigma,\sigma^*) \;=\; \bigl|\{(x,y)\in D(\sigma) : (y,x)\in D(\sigma^*)\}\bigr|,
\]
which counts direct machine-order disagreements between $\sigma$ and $\sigma^*$.
Note that $M(\sigma,\sigma^*)=0$ if and only if $D(\sigma)=D(\sigma^*)$, i.e., $\sigma=\sigma^*$.

\medskip
\noindent\textit{Key lemma.}
If $\lambda$ is non-optimal, then there exists an extreme critical arc $(x,y)\in D(\lambda)$ with $(y,x)\in D(\sigma^*)$.

\noindent\textit{Proof.}
Since $\lambda$ is non-optimal and every ranking $R$ refines $\le_2$, we have
$\mathbf{C}_{\max}(\lambda)\not\le_2 \mathbf{C}_{\max}(\sigma^*)$
(otherwise $\mathbf{C}_{\max}(\lambda)\le_R\mathbf{C}_{\max}(\sigma^*)$, contradicting non-optimality).
Thus either $C_{\max}^-(\lambda)>C_{\max}^-(\sigma^*)$ or $C_{\max}^+(\lambda)>C_{\max}^+(\sigma^*)$;
assume without loss of generality the former.
Let $P$ be a critical path in $G^-(\lambda)$, of length $C_{\max}^-(\lambda)$.
If every disjunctive arc $(u,v)$ of $P$ satisfied $(u,v)\in D(\sigma^*)$, then $P$ would be a valid directed path in $G^-(\sigma^*)$ of length $C_{\max}^-(\lambda)>C_{\max}^-(\sigma^*)$, contradicting the definition of $C_{\max}^-(\sigma^*)$.
Hence $P$ contains a disjunctive arc $(x,y)\in D(\lambda)\setminus D(\sigma^*)$.
Since for every machine pair exactly one of $(x,y),(y,x)$ belongs to $D(\sigma^*)$, we get $(y,x)\in D(\sigma^*)$.
Moreover $(x,y)$ is extreme critical because it lies on the critical path $P$ of $G^-(\lambda)$. \hfill$\square_{\text{lemma}}$

\medskip
\noindent\textit{Construction and termination.}
Build the sequence $\lambda_0=\sigma$, and at each step $k$:
\begin{itemize}
\item[(a)] If $\lambda_k$ is globally optimal, stop: a globally optimal order has been reached.
\item[(b)] Otherwise, apply the key lemma to obtain an extreme critical $(x,y)\in D(\lambda_k)$ with $(y,x)\in D(\sigma^*)$.
Set $\lambda_{k+1}=(\lambda_k)_{(x,y)}\in H(\lambda_k)$ (feasible by Theorem~\ref{thm1}).
\end{itemize}
After step (b), the inversion count satisfies $M(\lambda_{k+1},\sigma^*)=M(\lambda_k,\sigma^*)-1$:
the reversal removes $(x,y)$ from $D(\lambda_{k+1})$ (eliminating one inversion, since $(y,x)\in D(\sigma^*)$),
adds $(y,x)$ (which agrees with $\sigma^*$, adding no inversion),
and leaves all other arcs unchanged.

Since $M\ge 0$ and decreases by exactly one at each non-optimal step, within at most $M(\sigma,\sigma^*)$ steps the sequence must reach case (a): either a globally optimal $\lambda_k\ne\sigma^*$,
or $M(\lambda_k,\sigma^*)=0$ implying $\lambda_k=\sigma^*$.
In either case a globally optimal order is reached by a finite sequence of moves in $H$.
\end{proof}


\section{Neighbourhood Variants for IJSP}
\label{sec5}

We instantiate five concrete neighbourhood variants within the $H(\sigma)$ framework. They differ only in which subset of extreme critical arcs each one considers; all five inherit the feasibility, connectivity and no-loss-of-improving-neighbours guarantees of Section~\ref{sec4}. Figure~\ref{fig2} illustrates a critical block and the arcs included by each variant.

\begin{figure}[t]
\centering
\figorplaceholder{fig2_block_structure}{0.85\linewidth}{4cm}
\caption{Critical block structure and arcs considered by each neighbourhood. Arcs marked $\circ$ are unconditionally included; arcs marked with a $?$ are subject to a slack-based filter.}
\label{fig2}
\end{figure}

\subsection{$N_1$: All Extreme Critical Arcs}
\label{sec5.1}

$N_1$ includes every disjunctive arc $(x,y)$ such that $x$ immediately precedes $y$ on the same machine in $G(\sigma)$ \emph{and} the arc lies on a critical path of $G^-(\sigma)$ or $G^+(\sigma)$. No positional filter is applied within the block: both boundary and interior arcs of every extreme critical block are candidates. Generation cost: $O(|E|)$ per extreme graph.

\subsection{$N_2$: Block-Boundary Filter}
\label{sec5.2}

$N_2$ restricts $N_1$ to arcs at the \emph{boundary} of each extreme critical block. A block $B = \{o_1,\ldots,o_k\}$ is a maximal contiguous subsequence of operations on the same machine that all lie on a single extreme critical path. $N_2$ retains the two boundary arcs $(o_1,o_2)$ and $(o_{k-1},o_k)$ and discards interior arcs $(o_i,o_{i+1})$, $2 \leq i \leq k-2$. This is the direct adaptation of \cite{Nowicki1996} to the interval setting: by Proposition~\ref{prop2} applied simultaneously to $G^-(\sigma)$ and $G^+(\sigma)$, the discarded interior arcs cannot carry improving moves under any of the four rankings, so $N_2$ inherits the same theoretical guarantees as $N_1$.

\subsection{$N_3$: Subsequence Reversals}
\label{sec5.3}

$N_3$ extends the candidate set with reversals of \emph{contiguous subsequences} within a critical block. For block $B = \{o_1,\ldots,o_k\}$, in addition to the single-arc swaps of $N_1$, $N_3$ considers every contiguous pair $(o_i, o_j)$ with $j > i+1$ as a multi-step reversal candidate. This generates a strictly larger neighbourhood than $N_1$.

\subsection{$N_{\rm ext}$: Extended Nowicki--Smutnicki Neighbourhood}
\label{sec5.4}

$N_{\rm ext}$ adapts the extended neighbourhood of \cite{Nowicki2005}. It starts from $N_2$ (boundary arcs unconditionally) and additionally considers interior block arcs that pass a heads-and-tails viability check, denoted \texttt{isViableSwap}:
\begin{itemize}
\item Boundary arcs $(o_1,o_2)$ and $(o_{k-1},o_k)$ are always included (as in $N_2$).
\item Interior arcs $(o_i, o_{i+1})$, $2 \leq i \leq k-2$, are included only if $\texttt{isViableSwap}(o_i, o_{i+1})$ returns true.
\end{itemize}
$\texttt{isViableSwap}(x,y)$ computes an optimistic lower-bound estimate of the makespan after swapping $x$ and $y$ by recalculating only the local heads and tails of $x$ and $y$, without propagating changes through the full graph. Because the estimate is a lower bound, discarding a swap when the estimate exceeds the incumbent makespan is theoretically safe.

\subsection{$N_8$: Boundary with Tightened Slack Filter}
\label{sec5.5}

$N_8$ is a new variant proposed in this work. It uses the same boundary-arc skeleton as $N_2$, but applies an additional acceptance test based on the \emph{combined slack} of the two endpoints in both extreme graphs. An arc $(x,y)$ is accepted iff $F_{x}^{-} + F_{y}^{-} + F_{x}^{+} + F_{y}^{+}$ does not exceed an instance-dependent threshold computed at search time. The intent is to favour moves whose endpoints are critical in both extreme graphs simultaneously --- a structural cue that the two scenarios disagree on the path identity. As shown in Sections~\ref{sec6.5} and~\ref{sec6.6}, $N_8$ is competitive on the deterministic CRISP setting (where only one graph exists) but underperforms its peers in the genuinely interval-valued case.


\section{Experimental Study}
\label{sec6}

The experimental study is organised in three phases that address distinct questions.
\begin{description}
\item[Phase A (Sections~\ref{sec6.2} to~\ref{sec6.4}).] At a \emph{common} hyperparameter setting, we compare the four ranking operators on four neighbourhoods ($4 \times 4$ design), focusing on solution quality, the width of the returned interval, and the structural mechanism behind the LEX2--width effect.
\item[Phase B (Section~\ref{sec6.5}).] After an irace-based tuning of the hill-climbing local search separately for each of the five neighbourhoods, we compare the five neighbourhoods on equal footing. The tuning isolates the structural effect of the neighbourhood from confounding parameter choices.
\item[Phase C (Section~\ref{sec6.6}).] We validate Phase B on the 82 corresponding deterministic instances (CRISP), measuring the relative percentage deviation (RPD) from the best known solutions and comparing the relative ranking of the five neighbourhoods.
\end{description}

\subsection{Setup Common to All Phases}
\label{sec6.1}

\paragraph{Instances.} The benchmark suite contains 82 instances derived from classical JSP benchmarks (abz, ft, la and tai families \cite{Taillard1993,Beasley1990}) by perturbing each nominal processing time $p_{ij}^{*}$ with a $\pm 7.5\%$ uniform symmetric interval, $\mathbf{p}_{ij} = [0.925\,p_{ij}^{*},\, 1.075\,p_{ij}^{*}]$ (15\% relative width). Sizes range from $10 \times 5$ (ft10) to $50 \times 20$ (the largest tai instances, 1000 operations).

\paragraph{Algorithm.} Each configuration runs an iterated local search (ILS) with a hill-climbing inner loop, best-improvement neighbour selection, random block-reversal perturbation upon stagnation, and a stopping criterion based on a fixed number of generations without improvement. Thirty independent runs are executed for every (instance, configuration) pair.

\paragraph{Statistics.} Comparisons use the Friedman test (non-parametric repeated-measures ANOVA) followed by pairwise Wilcoxon signed-rank tests with Holm--Bonferroni correction. Effect size is reported as $r = |Z| / \sqrt{N}$ with the conventional thresholds (small $\geq 0.1$, medium $\geq 0.3$, large $\geq 0.5$). Each (instance, run) pair forms a paired block. The metric is the midpoint of the makespan interval, $(C^-_{\max} + C^+_{\max})/2$, which is a neutral common ground for all four operators.

\subsection{Phase A: Ranking Operator Comparison}
\label{sec6.2}

At a common hyperparameter setting, the four operators are compared on four neighbourhoods ($N_1$, $N_2$, $N_3$, $N_{\rm ext}$), $4 \times 4 = 16$ configurations, $16 \times 82 \times 30 = 39{,}360$ runs in total. Table~\ref{tab1} reports the global mean midpoint over all 82 instances. The Friedman test across operators yields $\chi^{2}(3) = 694.96,\, p < 10^{-150}$.

\begin{table}[t]
\centering
\caption{Phase A: global mean midpoint by neighbourhood and ranking operator (averaged over 82 instances $\times$ 30 runs, common hyperparameter setting). Lower is better.}
\label{tab1}
\begin{tabular}{l c c c c}
\hline
Neighbourhood & EV & LEX1 & LEX2 & YX \\
\hline
$N_1$         & 1888.9 & 1894.6 & 1889.5 & 1887.0 \\
$N_2$         & 1888.7 & 1895.2 & 1889.2 & 1887.1 \\
$N_3$         & 1907.3 & 1914.2 & 1908.5 & 1904.9 \\
$N_{\rm ext}$ & 1889.0 & 1891.0 & 1895.5 & 1887.2 \\
\hline
\end{tabular}
\end{table}

Pairwise tests on operators (Table~\ref{tab2}) show that LEX1 is significantly worse than the other three (medium-to-large effects), while EV and LEX2 are statistically equivalent ($p = 0.197$, $r = 0.026$). YX is the best on average, but its advantage over EV and LEX2 carries only small effects.

\begin{table}[t]
\centering
\caption{Phase A: pairwise Wilcoxon tests for ranking operators ($n = 2460$ paired blocks, Holm--Bonferroni corrected).}
\label{tab2}
\begin{tabular}{l l c c c c c l}
\hline
A & B & Mean A & Mean B & Diff & $p_{\rm adj}$ & $r$ & Magnitude \\
\hline
LEX1 & YX   & 1909.5 & 1901.5 & $+8.0$ & $<10^{-3}$ & 0.509 & large \\
EV   & LEX1 & 1903.3 & 1909.5 & $-6.2$ & $<10^{-3}$ & 0.415 & medium \\
LEX1 & LEX2 & 1909.5 & 1903.9 & $+5.6$ & $<10^{-3}$ & 0.349 & medium \\
LEX2 & YX   & 1903.9 & 1901.5 & $+2.4$ & $<10^{-3}$ & 0.159 & small \\
EV   & YX   & 1903.3 & 1901.5 & $+1.8$ & $<10^{-3}$ & 0.131 & small \\
EV   & LEX2 & 1903.3 & 1903.9 & $-0.6$ & $0.197$    & 0.026 & negligible \\
\hline
\end{tabular}
\end{table}

\subsection{Phase A: Interval Width by Operator}
\label{sec6.3}

The \emph{width} of the returned makespan interval, $C^+_{\max} - C^-_{\max}$, is a measure of solution \emph{uncertainty}: the narrower the interval, the more precisely the planner can predict the actual makespan, regardless of which scenario realises. Figure~\ref{fig3} visualises the distribution of width by operator across all neighbourhoods and instances; the operator effect is consistent: LEX2 produces the narrowest intervals.

\begin{figure}[t]
\centering
\includegraphics[width=0.6\linewidth]{fig4_width_by_operator}
\caption{Distribution of makespan interval width $C^+_{\max} - C^-_{\max}$ by ranking operator (Phase A; 5 neighbourhoods $\times$ 82 instances $\times$ best-of-30 = 1640 observations per operator). The dashed line marks the EV mean.}
\label{fig3}
\end{figure}

Table~\ref{tab3} quantifies the effect: LEX2 yields intervals $\sim 5\%$ narrower than EV and YX and $\sim 11\%$ narrower than LEX1. The ordering EV $\succ$ YX $\succ$ LEX1 $\succ$ LEX2 (in width) holds in every neighbourhood; the choice of neighbourhood has a negligible effect on width ($<3$ units within any operator).

\begin{table}[t]
\centering
\caption{Mean makespan interval width by ranking operator (Phase A; mean over 4 neighbourhoods $\times$ 82 instances).}
\label{tab3}
\begin{tabular}{l c c}
\hline
Operator & Mean width & vs.\ EV \\
\hline
EV   & 269.3 & --- \\
LEX1 & 286.4 & $+6.4\%$ \\
\textbf{LEX2} & \textbf{255.0} & $\mathbf{-5.3\%}$ \\
YX   & 267.5 & $-0.7\%$ \\
\hline
\end{tabular}
\end{table}

\subsection{Phase A: Why does LEX2 Narrow Intervals? A Float Analysis}
\label{sec6.4}

The natural hypothesis (H2) for the LEX2--width effect is that LEX2 leaves \emph{more structural slack} in $G^+(\sigma)$, which would let the worst-case path shorten. The alternative hypothesis (H1) is that LEX2 \emph{tightens} $G^+(\sigma)$ directly by accepting moves that reduce $C^+_{\max}$ even when $C^-_{\max}$ does not improve. To discriminate, we reconstruct $G^-(\sigma)$ and $G^+(\sigma)$ from the saved best-of-30 permutation for every (configuration, instance) pair and compute the float of every operation in both graphs.

Figure~\ref{fig4} reports two derived statistics: \emph{(a)} the mean number of critical operations in $G^+(\sigma)$ (i.e., operations with $F^{+} = 0$), and \emph{(b)} the mean per-operation float in $G^+(\sigma)$. The two panels point in the \emph{same} direction: LEX2 has \emph{more} critical operations in $G^+(\sigma)$ and \emph{less} per-operation slack. This is the opposite of what H2 predicts and exactly what H1 predicts.

\begin{figure}[t]
\centering
\includegraphics[width=0.95\linewidth]{fig5_float_analysis}
\caption{Structural float analysis in $G^+(\sigma)$ by ranking operator (Phase A; 1640 observations per operator). (a) Number of critical operations in $G^+(\sigma)$. (b) Mean per-operation float in $G^+(\sigma)$. LEX2 simultaneously produces \emph{more} critical operations and \emph{less} per-operation slack: the worst-case critical path is structurally tighter, not looser.}
\label{fig4}
\end{figure}

The asymmetry between the two extreme graphs is the structural fingerprint of LEX2: the operator pushes optimisation effort entirely toward the upper-bound graph $G^+(\sigma)$ (where it has more critical ops and less slack), at the cost of a relatively looser $G^-(\sigma)$ (fewer critical ops there). Both effects are mediated by the ranking, not by a coincidental property of the search trajectory.

\subsection{Phase B: Neighbourhood Comparison after Tuning}
\label{sec6.5}

Phase A used a common hyperparameter setting that prevents over-fitting to a single neighbourhood, but it leaves open the possibility that any of the five neighbourhoods would benefit from its own tuned configuration. To address this, we tune the hill-climbing parameters separately for each of the five neighbourhoods using irace \cite{Lopez2016}, fix the LEX2 ranking, and run the resulting tuned configurations on the 82 benchmark instances (30 runs per instance, $5 \times 82 \times 30 = 12{,}300$ observations).

Figure~\ref{fig5}(a) shows the per-instance RPD (relative percentage deviation from the best of all five neighbourhoods on the same instance and run): values of zero correspond to the per-instance winner, and higher values measure relative degradation. The Friedman test yields $\chi^{2}(4) = 1113.03,\, p < 10^{-150}$; mean ranks are $N_{\rm ext}$: 2.465, $N_1$: 2.667, $N_2$: 2.739, $N_8$: 3.479, $N_3$: 3.651. Pairwise tests (Table~\ref{tab4}) confirm a clean partition into a top group $\{N_{\rm ext}, N_1, N_2\}$, with $N_8$ and $N_3$ trailing by a medium-effect margin.

\begin{figure}[t]
\centering
\includegraphics[width=\linewidth]{fig5_combined}
\caption{Phase B (irace-tuned hill climbing). \textbf{(a)}~Per-instance RPD on the makespan midpoint by neighbourhood. Each box pools $82 \times 30 = 2{,}460$ observations; the per-instance baseline is the minimum midpoint across all five tuned neighbourhoods on the same run. Brackets indicate Wilcoxon significance; all pairs significant after Holm--Bonferroni correction (see Table~\ref{tab4} for effect sizes). \textbf{(b)}~Quality--runtime trade-off (N3 omitted for scale clarity): each point is the mean RPD vs.\ mean median runtime aggregated over all instance groups; error bars are $\pm 1$ SE. $N_{\rm ext}$ achieves the lowest RPD at moderate runtime; $N_8$ is the fastest but trails on quality by a medium effect.}
\label{fig5}
\end{figure}

\begin{table}[t]
\centering
\caption{Phase B: pairwise Wilcoxon tests for tuned neighbourhoods ($n = 2460$ paired blocks, Holm--Bonferroni corrected). All differences are statistically significant; the magnitude of the effect indicates the practical relevance.}
\label{tab4}
\begin{tabular}{l l c c c c c l}
\hline
A & B & Mean A & Mean B & Diff & $p_{\rm adj}$ & $r$ & Magnitude \\
\hline
$N_3$         & $N_{\rm ext}$ & 1898.8 & 1882.1 & $+16.7$ & $<10^{-3}$ & 0.530 & large \\
$N_1$         & $N_3$         & 1884.0 & 1898.8 & $-14.7$ & $<10^{-3}$ & 0.469 & medium \\
$N_{\rm ext}$ & $N_8$         & 1882.1 & 1896.9 & $-14.9$ & $<10^{-3}$ & 0.458 & medium \\
$N_1$         & $N_8$         & 1884.0 & 1896.9 & $-12.9$ & $<10^{-3}$ & 0.407 & medium \\
$N_2$         & $N_3$         & 1887.6 & 1898.8 & $-11.2$ & $<10^{-3}$ & 0.385 & medium \\
$N_2$         & $N_8$         & 1887.6 & 1896.9 & $-9.4$  & $<10^{-3}$ & 0.320 & medium \\
$N_2$         & $N_{\rm ext}$ & 1887.6 & 1882.1 & $+5.5$  & $<10^{-3}$ & 0.171 & small \\
$N_1$         & $N_{\rm ext}$ & 1884.0 & 1882.1 & $+2.0$  & $<10^{-3}$ & 0.097 & negligible \\
$N_1$         & $N_2$         & 1884.0 & 1887.6 & $-3.5$  & $<10^{-3}$ & 0.076 & negligible \\
$N_3$         & $N_8$         & 1898.8 & 1896.9 & $+1.8$  & $<10^{-3}$ & 0.067 & negligible \\
\hline
\end{tabular}
\end{table}

The picture that emerges is sharper than under the common hyperparameter setting of Phase A: $N_{\rm ext}$ overtakes $N_1$ and $N_2$ as the strongest neighbourhood after tuning (small but consistent effect), and the new $N_8$ variant is a medium-effect step \emph{worse} than the top group. Figure~\ref{fig5}(b) further shows that $N_{\rm ext}$ achieves the lowest RPD at moderate runtime ($\approx$34\,s median), while $N_8$ is the fastest ($\approx$15\,s) but pays the medium-effect quality penalty; $N_1$ sits between them in both dimensions. The latter is, at first sight, surprising: $N_8$'s structural cue --- arcs whose endpoints are critical in both extreme graphs --- seemed promising. Section~\ref{sec6.6} shows that the cue is indeed informative when only \emph{one} graph exists (CRISP), but it filters out moves that are valuable specifically because they target \emph{the disagreement} between $G^-(\sigma)$ and $G^+(\sigma)$.

\subsubsection{Convergence}
\label{sec6.5.1}

Figure~\ref{fig6} compares the per-generation convergence profiles of the five tuned neighbourhoods, normalised so that every trace starts at $1.000$ at generation $0$. The top group ($N_{\rm ext}, N_1, N_2$) is essentially indistinguishable in convergence speed, all reaching their plateau by generation $\sim 50$. $N_3$ converges slightly slower and to a worse plateau. $N_8$ exhibits a clear gap: it converges noticeably more slowly and stalls $\sim 4\%$ above the top group. This is consistent with its weaker structural acceptance test: many improving moves are rejected by the slack threshold, leading to longer search trajectories at lower yield.

\begin{figure}[t]
\centering
\includegraphics[width=0.95\linewidth]{fig6_convergence}
\caption{Per-generation normalised best-makespan profiles, by neighbourhood (Phase B; mean over 82 instances). The shaded ribbons represent $\pm 1$ standard error of the mean. The top group ($N_{\rm ext}, N_1, N_2$) overlaps tightly; $N_8$ converges noticeably slower than the top group; $N_3$ plateaus at a worse value.}
\label{fig6}
\end{figure}

\subsection{Phase C: Validation on Deterministic Instances (CRISP)}
\label{sec6.6}

To check whether the relative ranking of neighbourhoods carries over to the deterministic setting, we apply the same five tuned configurations of Phase B to the 82 CRISP counterparts of the benchmark instances ($\mathbf{p}_{ij} = [p_{ij}^{*}, p_{ij}^{*}]$, equivalently $G^-(\sigma) = G^+(\sigma)$). The metric is RPD against the best known solutions from \cite{Taillard1993,Beasley1990,TaillardBKS}:
\begin{equation}
{\rm RPD}(\sigma) = \frac{C_{\max}(\sigma) - {\rm BKS}}{{\rm BKS}} \times 100.
\end{equation}

Table~\ref{tab5} reports global summary statistics. The relative ordering of the five neighbourhoods on CRISP is $N_{\rm ext} \succeq N_8 \succeq N_1 \succeq N_2 \succeq N_3$, with $N_{\rm ext}$ leading on mean RPD (3.18\%) and $N_8$ leading on the number of optimal solutions (6 vs.\ 4 for the others). $N_8$ is genuinely competitive on CRISP --- a marked contrast with its Phase B performance.

\begin{table}[t]
\centering
\caption{Phase C: CRISP validation. RPD against best known solutions over the 82 instances and 30 runs.}
\label{tab5}
\begin{tabular}{l c c c c c}
\hline
Neighbourhood & Mean RPD\,(\%) & Median RPD\,(\%) & SD & Max\,(\%) & \# optimal \\
\hline
$N_{\rm ext}$ & 3.175 & 2.988 & 2.195 & 9.101 & 4 \\
$N_8$         & 3.290 & 2.960 & 2.372 & 9.751 & \textbf{6} \\
$N_1$         & 3.359 & 3.219 & 2.336 & 9.158 & 4 \\
$N_2$         & 3.374 & 3.197 & 2.307 & 9.751 & 4 \\
$N_3$         & 3.802 & 3.729 & 2.523 & 10.130 & 4 \\
\hline
\end{tabular}
\end{table}

Figure~\ref{fig7} resolves the global picture by instance group. Three observations stand out:
\begin{enumerate}
\item The top three neighbourhoods ($N_{\rm ext}$, $N_1$, $N_2$) maintain the same partition observed in Phase B, with $N_{\rm ext}$ leading on most groups.
\item $N_8$ is best or co-best on multiple small/medium-sized groups (la21, la38, la40, tai15$\times$15, tai20$\times$15, tai20$\times$20) and matches $N_{\rm ext}$ on tai30$\times$15 and abz8 --- precisely those groups where the $N_2$ skeleton already does well.
\item Difficulty grows with size: tai30$\times$20 is the hardest group across the board (RPD 6.8--8.2\%), with all five neighbourhoods converging to similar values.
\end{enumerate}

\begin{figure}[t]
\centering
\includegraphics[width=\linewidth]{fig7_crisp_rpd_heatmap}
\caption{Phase C: CRISP RPD\,(\%) against BKS, by instance group and neighbourhood. Cells display mean RPD across the instances of each group; colours saturate the green--red gradient at $0$ (BKS) and $12\%$ respectively. $N_8$ is competitive when only one extreme graph exists.}
\label{fig7}
\end{figure}

The CRISP validation supports an interpretation of the $N_8$ behaviour: its slack-based acceptance filter is a sound structural cue when there is a single graph (deterministic), but \emph{filters out exactly the moves that are most informative in the interval setting} --- those that resolve disagreement between $G^-(\sigma)$ and $G^+(\sigma)$. This is, to the best of our knowledge, the first experimental evidence that a CRISP-good neighbourhood can be IJSP-bad for purely structural reasons.


\section{Discussion}
\label{sec7}

\subsection{Practical Recommendations}
\label{sec7.1}

The experimental results lead to a small set of actionable recommendations:
\begin{enumerate}
\item \textbf{Use $N_{\rm ext}$ if compute budget allows; $N_2$ if not.} After per-neighbourhood tuning, $N_{\rm ext}$ is the strongest variant on intervals (and on CRISP), but the $N_{\rm ext}$--$N_1$ gap is small ($r = 0.097$) and the $N_{\rm ext}$--$N_2$ gap is small ($r = 0.171$). $N_2$ remains attractive for its lighter generation cost.
\item \textbf{Use LEX2 if interval width is the goal, YX if midpoint is the goal.} The two objectives are not aligned: LEX2 wins on width with negligible cost on midpoint (vs.\ EV); YX wins on midpoint with a small cost on width (vs.\ LEX2).
\item \textbf{Avoid $N_8$ for genuinely interval-valued problems.} Its slack-based filter is informative on CRISP but counter-productive on intervals --- it filters out disagreement-resolving moves.
\item \textbf{Avoid $N_3$ in either setting.} The multi-step subsequence reversals it adds beyond $N_1$ are dominated empirically: more candidates, no improvement.
\end{enumerate}

\subsection{Theoretical Connection: Why LEX2 Narrows Intervals}
\label{sec7.2}

Proposition~\ref{prop2} guarantees that reversing a non-extreme-critical arc cannot improve either bound. LEX2 specifically targets $C^+_{\max}$, accepting moves that reduce $C^+_{\max}$ even when $C^-_{\max}$ stays put. Over many local search steps, this creates a bias toward configurations where $G^+(\sigma)$ has a shorter critical path --- and, by the structure of interval arithmetic, a tighter bound on the overall interval width. The float analysis of Section~\ref{sec6.4} provides direct empirical confirmation: LEX2 solutions have more critical operations in $G^+(\sigma)$ and less per-operation slack, while $G^-(\sigma)$ becomes relatively looser. The asymmetry between the two extreme graphs is a structural fingerprint of the LEX2 ranking rule.

\subsection{Why $N_3$ Underperforms}
\label{sec7.3}

$N_3$ is counter-intuitive: it is a strict superset of $N_1$ in the move-generation sense (single-arc swaps plus multi-step subsequence reversals), yet it produces worse solutions. The \texttt{isViableSwap} heuristic is calibrated for single-arc swaps; for multi-arc reversals it becomes excessively optimistic, allowing many spurious candidates through into the evaluated pool. Combined with a best-improvement search strategy, a larger pool diluted with poor candidates reduces the probability of selecting a genuinely improving move quickly, increasing wasted evaluations per step.

\subsection{Why $N_8$ Is CRISP-Good but IJSP-Bad}
\label{sec7.4}

The slack-based filter of $N_8$ accepts an arc only when the four endpoint floats are jointly small. In the deterministic setting, this is a sound structural cue: only one extreme graph exists, and the filter selects moves that simultaneously affect the (single) critical path at both endpoints. In the interval setting, however, the most informative moves are precisely those that reduce $C^+_{\max}$ without affecting $C^-_{\max}$ (or vice versa) --- i.e., moves whose endpoints are critical in only \emph{one} of the two extreme graphs. The $N_8$ filter rejects these moves, costing the search the very degrees of freedom that distinguish the IJSP from the JSP.

\subsection{Limitations and Future Work}
\label{sec7.5}

\begin{itemize}
\item \textbf{Larger Taillard instances.} The current study covers up to $\text{tai}50 \times 20$ (1000 operations); $\text{tai}100 \times 20$ (2000 operations) is in progress.
\item \textbf{Other uncertainty levels.} The 15\% uniform interval width is a single noise level; sensitivity to other noise widths is left open.
\item \textbf{Hybrid algorithms.} The combinations $({\rm LEX2}, N_{\rm ext})$ and $({\rm YX}, N_2)$ are natural starting points for a dedicated uncertainty-aware metaheuristic.
\item \textbf{Other rankings.} Additional admissible orders on intervals \cite{Bustince2013} could be tested with the same protocol; the structural mechanism uncovered for LEX2 should generalise to any operator that strictly prioritises one extreme.
\end{itemize}


\section{Conclusions}
\label{sec8}

We have presented a theoretical framework for neighbourhood-based local search in the Interval Job Shop Scheduling Problem grounded in the concept of \emph{extreme critical paths} --- arcs that are critical in either the best-case or worst-case scenario graph. The resulting neighbourhood $H(\sigma)$ satisfies feasibility, connectivity and the no-loss-of-improving-neighbours property for all four admissible interval rankings considered (EV, LEX1, LEX2, YX) simultaneously.

Five neighbourhood variants ($N_1$, $N_2$, $N_3$, $N_{\rm ext}$, $N_8$) and four ranking operators were evaluated on 82 benchmark instances of up to 1000 operations, in three experimental phases: a $4 \times 4$ comparison at a common parameter setting, a five-way comparison after irace-based tuning, and a CRISP validation on the corresponding deterministic instances. The main findings are:
\begin{enumerate}
\item After tuning, $N_{\rm ext}$ is the strongest neighbourhood, followed by $N_1$ and $N_2$ within a small effect; $N_8$ and $N_3$ trail by a medium effect.
\item LEX2 systematically yields makespan intervals 5--11\% narrower than the alternatives, with negligible cost on the midpoint. A float analysis of both extreme graphs explains this mechanistically: LEX2 directly tightens the worst-case critical path rather than leaving more structural slack.
\item The new $N_8$ variant is competitive on CRISP but underperforms on intervals --- a counter-intuitive finding that traces back to its slack-based filter rejecting precisely the moves that resolve disagreement between $G^-(\sigma)$ and $G^+(\sigma)$.
\end{enumerate}
These results provide actionable guidance for practitioners implementing local search for uncertain job shop scheduling: the choice of ranking operator can be used not only to optimise expected makespan, but also to directly control solution uncertainty, and the choice of neighbourhood should account for whether the problem is genuinely interval-valued or only nominally uncertain.


\section*{CRediT authorship contribution statement}

\textbf{Hern\'an D\'{\i}az Rodr\'{\i}guez}: Conceptualization, Methodology, Software, Formal analysis, Investigation, Data curation, Visualization, Writing -- original draft, Writing -- review and editing.


\section*{Declaration of competing interest}

The author declares that there are no known competing financial interests or personal relationships that could have appeared to influence the work reported in this paper.


\section*{Data availability}

[TODO: Confirm or replace.] The benchmark instances, source code of the local-search engine, irace tuning scripts, and the raw experimental data underlying every figure and table in this article are openly available at \texttt{https://github.com/<TODO-org>/<TODO-repo>}. The classical JSP benchmarks used to derive the interval instances are taken from the OR-Library and from \cite{TaillardBKS}.


\section*{Acknowledgements}

[TODO: Funding sources, project IDs, computing resources.]


%% If you have bib database file and want bibtex to generate the
%% bibitems, please use
%%
%%   \bibliographystyle{elsarticle-num}
%%   \bibliography{<your bibdatabase>}

\bibliographystyle{elsarticle-num}
\bibliography{references}

\end{document}
